% !PRINTED 341
\weberpage{341}
%<<CONT
folgt, daß $a_{11}$ durch $a_{22}$ teilbar ist, und wir setzen $a_{11} = a_{22}a$.
Dann haben wir die Basis von $\mathfrak{a}$:
\weq{8}{\begin{aligned} \alpha_1 &= a_{22}a,\\ \alpha_2 &= a_{12} + a_{22}\theta. \end{aligned}}
Es ist dann nach Bd.\,II, §\,164, (12):
\weq{9}{N(\mathfrak{a}) = a_{22}^2\,a,}
und wenn wir für $(1, \theta)$ die Basis §\,90, (9) nehmen, so ist
\[
\theta + \theta' = 0 \text{ oder } = -1,
\]
\weq{10}{\theta\theta' = -\tfrac14\Delta \text{ oder } = \frac{1-\Delta}{4}.}
Je nachdem $\Delta \equiv 0$, oder $\equiv 1 \pmod 4$ ist, und
\[
(\theta - \theta')^2 = \Delta.
\]
Es wird dann in diesen beiden Fällen, wenn man (8) in (3)
substituiert und (9) berücksichtigt:
\[
a_{22}b = 2a_{12} \text{ oder } = 2a_{12} - a_{22},
\]
und folglich in beiden Fällen:
\weq{11}{\alpha_1 = a_{22}a, \quad \alpha_2 = a_{22}\,\frac{b+\sqrt{\Delta}}{2},}
und die Form
\[
T = (a,\,b,\,c)
\]
hat hier einen \gesperrt{positiven ersten Koeffizienten}.

Setzen wir
\weq{12}{\omega = -\frac{\alpha_2}{\alpha_1} = -\frac{b+\sqrt{\Delta}}{2a},}
so ist $\omega$ Wurzel der quadratischen Gleichung:
\weq{13}{a\omega^2 + b\omega + c = 0}
und
\weq{14}{\lambda = a_{22}a(x-\omega y)}
ist ein dem Ideal $\mathfrak{a}$ entsprechendes Funktional und zugleich eine
Basisform von $\mathfrak{a}$. $a$ ist die kleinste positive ganze Zahl, für die
$a\omega$ eine ganze Zahl wird.

Die notwendige und hinreichende Bedingung dafür, daß zwei
ganze Zahlen $\alpha_1, \alpha_2$ des Körpers $\Omega$ die Basis eines Ideals bilden,
besteht darin, daß zunächst $\alpha_1, \alpha_2$ eine Basis des Körpers $\Omega$ ist,
und daß, wenn $\omega_1, \omega_2$ eine Minimalbasis dieses Körpers ist,
\[
\alpha_1\omega_1,\quad \alpha_2\omega_1,\quad \alpha_1\omega_2,\quad \alpha_2\omega_2
\]
durch die Linearform
\[
\lambda = \alpha_1 t_1 + \alpha_2 t_2
\]
%CONT>>
